\clearpage
\section*{Resource sharing plan}

\instructions{Investigators seeking \$500,000 or more in direct costs
 (exclusive of consortium F\&A) in any budget period are expected to
 include a brief 1-paragraph description of how final research data
 will be shared, or explain why data-sharing is not possible (for
 example human subject concerns, the Small Business Innovation
 Development Act provisions, etc.). Specific FOAs may require that
 all applications include this information regardless of the dollar
 level. Applicants are encouraged to read the FOA carefully and
 discuss their data-sharing plan with their program contact at the
 time they negotiate an agreement with the Institute/Center (IC)
 staff to accept assignment of their application.}

\subsection*{Data sharing plan}

\subsection*{Software dissemination plan}

\instructions{A software dissemination plan, with appropriate
 timelines, must be included in the application. There is no
 prescribed single license for software produced through grants
 responding to this announcement. However, NIH does have goals for
 software dissemination, and reviewers will be instructed to evaluate
 the dissemination plan relative to these goals:
\begin{enumerate}
\item The software should be freely available to biomedical
 researchers and educators in the non-profit sector, such as
 institutions of education, research institutions, and government
 laboratories.
\item The terms of software availability should permit the
 dissemination and commercialization of enhanced or customized
 versions of the software, or incorporation of the software or pieces
 of it into other software packages.
\item To preserve utility to the community, the software should be
 transferable such that another individual or team can continue
 development in the event that the original investigators are
 unwilling or unable to do so.
\item The terms of software availability should include the ability of
 researchers to modify the source code and to share modifications
 with other colleagues. An applicant should take responsibility for
 creating the original and subsequent ``official'' versions of a
 piece of software.
\item To further enhance the potential impact of their software,
 applicants may consider proposing a plan to manage and disseminate
 the improvements or customizations of their tools and resources by
 others. This proposal may include a plan to incorporate the
 enhancements into the ``official'' core software, may involve the
 creation of an infrastructure for plug-ins, or may describe some
 other solution.
\end{enumerate}
The adequacy of the software sharing plans will be considered by
Program staff when making recommendations about funding applications.
In making such considerations, prior to funding, program staff may
negotiate modifications of software sharing plans with the Principal
Investigator before recommending funding of an application. Any
software dissemination plans represent a commitment by the institution
(and its subcontractors as applicable) to support and abide by the
plan. The final version of any accepted software sharing plans will
become a condition of the award of the grant. The effectiveness of
software sharing may be evaluated as part of the administrative review
of each Non-Competing Grant Progress Report (PHS 2590).}
